\setlength{\emergencystretch}{3em}

% ============================================================
%  CHAPTER 4 — SPRINT 2
% ============================================================

\chapter[Chapter 4: Sprint 2 — AI Scan \& Diagnosis]{Chapter 4: Sprint 2 — AI Scan\\[12pt]
	\& Diagnosis}
\begin{center}
	\includegraphics[width=0.85\textwidth, keepaspectratio]{chap4_cover.jpg}
\end{center}

\newpage

% ============================================================
\section{Introduction}
% ============================================================
This chapter presents the second development sprint, focusing on the core AI-powered skin lesion scanner. The diagnostic pipeline leverages the Xception convolutional neural network, trained on the HAM10000 dataset and enhanced with Grad-CAM explainability. The integration tasks, frontend scanning interface, and detailed AI model evaluations are documented in the following sections.

% ============================================================
\section{Objectives}
% ============================================================
The main objectives to be accomplished during this sprint are:
\begin{itemize}
	\item Develop an intuitive frontend interface allowing users to upload skin lesion images.
	\item Implement an interactive body map for users to optionally pinpoint the location of the lesion, providing contextual data for the scan.
	\item Integrate the trained Xception deep learning model to return a primary and secondary diagnosis with confidence scores.
	\item Implement Grad-CAM (Gradient-weighted Class Activation Mapping) to generate visual heatmaps, explaining exactly which regions of the image the AI analyzed.
	\item Automatically store the scan results and contextual data in the user's history.
\end{itemize}

% ============================================================
\section{Design}
% ============================================================
This section details the functional and static design of the core scanning process.

% ────────────────────────────────────────────────────────────
\subsection{Use-Cases Modeling}
% ────────────────────────────────────────────────────────────
Figure~\ref{fig:sprint2_usecase} illustrates the use-cases diagram for the scanning process.

\begin{figure}[H]
	\centering
	\optionalimage[0.70\textwidth]{chap_4/images/sp2 use case.png}{Sprint 2 Use-Cases Diagram}
	\caption{\textit{Use-Cases Diagram for Sprint 2}}
	\label{fig:sprint2_usecase}
\end{figure}

% ············································
\subsubsection{Perform Skin Lesion Scan}
% ············································
\begin{table}[H]
	\centering
	\footnotesize
	\renewcommand{\arraystretch}{1.5}
	\begin{tabular}{|p{3cm}|p{11.5cm}|}
		\hline
		\rowcolor[HTML]{1F3864}
		{\color{white}\textbf{Use Case}} & {\color{white}\textbf{Perform Skin Lesion Scan}} \\
		\hline
		\textbf{Actor} & User \\
		\hline
		\textbf{Precondition} & The user is authenticated and on the Scan page. \\
		\hline
		\textbf{Postcondition} & The AI returns a diagnosis, a Grad-CAM heatmap, and saves the record. \\
		\hline
		\textbf{Nominal Scenario} & 
		1. The user uploads an image of a skin lesion.\newline
		2. (Optional) The user selects the affected area on the body map.\newline
		3. The system processes the image using the Deep Learning model.\newline
		4. The system generates a Grad-CAM heatmap.\newline
		5. The results (diagnosis, confidence, heatmap) are displayed.\newline
		6. The record is stored in the database history. \\
		\hline
		\textbf{Alternative} & 
		If the uploaded file is corrupted or not a valid image format, the system rejects it and displays an error alert. \\
		\hline
	\end{tabular}
	\caption{\textit{Textual Description: Perform Skin Lesion Scan}}
\end{table}

% ────────────────────────────────────────────────────────────
\subsection{Activity Diagram}
% ────────────────────────────────────────────────────────────
Figure~\ref{fig:sprint2_activity} presents the activity diagram modeling the scanning process.

\begin{figure}[H]
	\centering
	\optionalimage[0.82\textwidth]{chap_4/images/sp2 act diagram (1).png}{Activity Diagram for Scanning}
	\caption{\textit{Activity Diagram for Scanning Process}}
	\label{fig:sprint2_activity}
\end{figure}

% ────────────────────────────────────────────────────────────
\subsection{Class Diagram}
% ────────────────────────────────────────────────────────────
Figure~\ref{fig:sprint2_class} shows the conceptual class diagram for the entities managed during this sprint.

\begin{figure}[H]
	\centering
	\optionalimage[0.55\textwidth]{chap_4/images/sp2 class.png}{Sprint 2 Class Diagram}
	\caption{\textit{Conceptual Class Diagram for Sprint 2}}
	\label{fig:sprint2_class}
\end{figure}

% ────────────────────────────────────────────────────────────
\subsection{Sequence Diagram}
% ────────────────────────────────────────────────────────────
Figure~\ref{fig:sprint2_seq_scan} outlines the interactions between the React client, the Flask API, the Deep Learning Prediction Service, and the MongoDB database during an active scan request.

\begin{figure}[H]
	\centering
	\optionalimage[0.85\textwidth]{chap_4/images/seq sp2 (1).png}{Scanning Sequence Diagram}
	\caption{\textit{Sequence Diagram: Scanning and Prediction Flow}}
	\label{fig:sprint2_seq_scan}
\end{figure}

% ============================================================
\section{Tasks}
% ============================================================
The specific development tasks required to fulfill the objectives of Sprint 2 are defined in the sprint backlog below.

\begin{table}[H]
	\centering
	\small
	\renewcommand{\arraystretch}{1.2}
	\begin{tabular}{|>{\centering\arraybackslash}p{0.8cm}|>{\raggedright\arraybackslash}p{4.5cm}|>{\raggedright\arraybackslash}p{6.8cm}|}
		\hline
		\rowcolor[HTML]{1F3864}
		{\color{white}\textbf{ID}} & {\color{white}\textbf{User Story}} & {\color{white}\textbf{Task Description}} \\
		\hline
		\textbf{5} & As a user, I want to upload/capture a skin lesion image and optionally pinpoint its location on a body map. & Design the Image Upload component (drag/drop and camera capture). \\
		\cline{3-3}
		& & Develop the interactive SVG Body Map selector. \\
		\hline
		\textbf{6} & As a user, I want to receive an AI-powered preliminary diagnosis with confidence scores. & Develop the Flask prediction route (/upload-diagnosis). \\
		\cline{3-3}
		& & Integrate the pre-trained Xception model into the Flask backend. \\
		\hline
		\textbf{7} & As a user, I want to view a visual heatmap over my image to understand exactly which areas the AI analyzed. & Implement the Grad-CAM extraction script targeting the final convolutional layer. \\
		\cline{3-3}
		& & Develop the Results Display UI to render the diagnoses and base64 heatmap overlays. \\
		\hline
	\end{tabular}
	\caption{\textit{Sprint 2 Backlog}}
	\label{tab:sprint2_backlog}
\end{table}

% ============================================================
\section{Development}
% ============================================================
This section showcases the developed user interfaces for the scanning process and validates the backend logic through API testing.

% ────────────────────────────────────────────────────────────
\subsection{User Interfaces}
% ────────────────────────────────────────────────────────────

\paragraph{Scan and Body Map Interface}
The upload interface offers a clean drag-and-drop zone. Alongside the image upload, users interact with an SVG-based Body Map component to optionally select the general location of the lesion, which provides helpful medical context. Figure~\ref{fig:sp2_body} and Figure~\ref{fig:ui_scan} present the interactive body map and the scan interface, respectively.

\begin{figure}[H]
	\centering
	\optionalimage[0.50\textwidth]{chap_4/images/body.png}{Interactive Body Map}
	\caption{\textit{Interactive Body Map for Scan Location}}
	\label{fig:sp2_body}
\end{figure}

\begin{figure}[H]
	\centering
	\optionalimage[0.50\textwidth]{chap_4/images/upload_pic.png}{Scan Interface}
	\caption{\textit{Skin Lesion Image Upload Interface}}
	\label{fig:ui_scan}
\end{figure}

\paragraph{Results and Explanation Interface}
Once the prediction is processed, the system transitions to the Results Display. This interface prominently shows the uploaded image alongside the generated Grad-CAM heatmap, making the AI's focal points visually explicit. It also details the primary and secondary diagnoses with their respective confidence scores. Figure~\ref{fig:ui_results} shows this results interface.

\begin{figure}[H]
	\centering
	\optionalimage[0.85\textwidth]{chap_4/images/results.png}{Results Interface}
	\caption{\textit{Diagnostic Results Display}}
	\label{fig:ui_results}
\end{figure}

% ────────────────────────────────────────────────────────────
\subsection{API Testing}
% ────────────────────────────────────────────────────────────
The deep learning pipeline was verified using Postman by targeting the \texttt{/api/upload-diagnosis} endpoint.

\paragraph{Inference Endpoint Validation}
Uploading a valid dermoscopic image via a \texttt{multipart/form-data} POST request successfully triggers the prediction service. The endpoint processes the image through the Xception model and returns a JSON payload containing the \texttt{diagnosis}, \texttt{confidence}, \texttt{diagnosis\_2}, and a \texttt{grad\_cam\_image} formatted as a Base64 string (Status 200 OK). Figure~\ref{fig:api_scan} shows the Postman response from the inference endpoint.

\begin{figure}[H]
	\centering
	\optionalimage[0.85\textwidth]{chap_4/images/api_scan.png}{Postman Scan Test}
	\caption{\textit{Postman Test: Deep Learning Inference API}}
	\label{fig:api_scan}
\end{figure}

% ============================================================
\section{AI Model Development}
% ============================================================
The intelligence of the DermaAI platform relies on a highly optimized Convolutional
Neural Network pipeline. This section provides an in-depth account of all decisions
made during the design, training, and evaluation phases, grounded in the actual
experimental results obtained from the training notebook~\cite{HAM10000, TensorFlow}.

% ────────────────────────────────────────────────────────────
\subsection{Dataset: HAM10000}
% ────────────────────────────────────────────────────────────
The model was trained on the \textbf{HAM10000} (Human Against Machine with 10,000
training images) dataset~\cite{HAM10000}, the internationally recognized benchmark
for dermatological AI research. It contains \textbf{10,015} dermoscopic images across
\textbf{7 clinically relevant diagnostic classes}. The dataset was partitioned using
stratified sampling to preserve class distribution ratios across all splits, as
summarized in Table~\ref{tab:dataset_split}.

\begin{table}[H]
	\centering
	\renewcommand{\arraystretch}{1.5}
	\begin{tabular}{|l|c|c|l|}
		\hline
		\rowcolor[HTML]{1F3864}
		{\color{white}\textbf{Split}} &
		{\color{white}\textbf{Images}} &
		{\color{white}\textbf{Ratio}} &
		{\color{white}\textbf{Purpose}} \\
		\hline
		Training   & 8,012 & 80\% & Model weight optimization \\
		\hline
		Validation & 1,001 & 10\% & Hyperparameter tuning \& early stopping \\
		\hline
		Test       & 1,002 & 10\% & Final unbiased evaluation \\
		\hline
		\textbf{Total} & \textbf{10,015} & \textbf{100\%} & \\
		\hline
	\end{tabular}
	\caption{\textit{HAM10000 Dataset Split Distribution}}
	\label{tab:dataset_split}
\end{table}

The class distribution within the raw training set is severely imbalanced, as shown
in Table~\ref{tab:class_dist}. The dominant class (NV --- Melanocytic Nevi) accounts
for 66.95\% of training samples, while rare classes such as Dermatofibroma (DF) and
Vascular Lesions (VASC) represent less than 2\% each. A naive model trained on this
distribution would be strongly biased toward the majority class, producing clinically
unreliable results for rare but dangerous conditions.

\begin{table}[H]
	\centering
	\renewcommand{\arraystretch}{1.5}
	\begin{tabular}{|l|l|c|c|}
		\hline
		\rowcolor[HTML]{1F3864}
		{\color{white}\textbf{Code}} &
		{\color{white}\textbf{Class Name}} &
		{\color{white}\textbf{Train Samples}} &
		{\color{white}\textbf{\% of Training Set}} \\
		\hline
		NV    & Melanocytic Nevi (Benign mole)      & 5,364 & 66.95\% \\
		\hline
		MEL   & Melanoma (Malignant)                &   890 & 11.11\% \\
		\hline
		BKL   & Benign Keratosis-like Lesions       &   879 & 10.97\% \\
		\hline
		BCC   & Basal Cell Carcinoma                &   411 &  5.13\% \\
		\hline
		AKIEC & Actinic Keratoses                   &   262 &  3.27\% \\
		\hline
		VASC  & Vascular Lesions                    &   114 &  1.42\% \\
		\hline
		DF    & Dermatofibroma                      &    92 &  1.15\% \\
		\hline
	\end{tabular}
	\caption{\textit{Raw Class Distribution in the Training Set}}
	\label{tab:class_dist}
\end{table}

Figure~\ref{fig:dataset_samples} presents representative dermoscopic images from
each of the seven diagnostic classes, illustrating the high degree of visual
similarity between several malignant and benign lesion types that makes this
classification task clinically challenging.

\begin{figure}[H]
	\centering
	\optionalimage[0.85\textwidth]{chap_4/images/dataset_samples.png}{Representative Dataset Samples}
	\caption{\textit{Representative Samples for the Seven Diagnostic Classes}}
	\label{fig:dataset_samples}
\end{figure}

The sample distribution across each of the seven diagnostic classes illustrates the severe class imbalance and the high dominance of the NV class.

% ────────────────────────────────────────────────────────────
\subsection{Class Imbalance Strategy}
% ────────────────────────────────────────────────────────────
To counteract the severe class imbalance, a data-level oversampling strategy was
implemented via a custom \texttt{balance()} function. A target of \textbf{1,006
	samples per class} was established. For all minority classes, the function uses
Keras \texttt{ImageDataGenerator} to synthesize augmented copies using the
following transformations:

\begin{itemize}
	\item Horizontal flips.
	\item Random rotations up to $20^{\circ}$.
	\item Width and height shifts up to 20\%.
	\item Zoom augmentation up to 20\%.
\end{itemize}

After balancing, the effective training set grew from 8,012 raw images to
\textbf{7,042 balanced images} (7 classes $\times$ 1,006 samples each). All images
were resized to a fixed resolution of \textbf{480 $\times$ 480 pixels} with 3 RGB
channels, providing sufficient spatial detail for the depthwise separable convolutions
to extract fine-grained skin texture features.

% ────────────────────────────────────────────────────────────
\subsection{Architecture Selection: Xception Model}
% ────────────────────────────────────────────────────────────
The architecture selection was made following a comparative evaluation of leading CNN models available in the Keras Applications module (such as VGG16, ResNet50, and EfficientNet). \textbf{Xception}~\cite{Xception} was ultimately selected for the following reasons:
\begin{itemize}
	\item \textbf{Depthwise Separable Convolutions:} Xception decomposes standard
	convolutions into separate depthwise and pointwise steps, reducing
	parameter count while improving representational capacity --- critical for
	distinguishing subtle texture differences between malignant and benign lesions.
	\item \textbf{ImageNet Pre-trained Weights:} The backbone was initialized with
	ImageNet weights, providing rich generic feature representations (edges,
	gradients, textures) that transfer effectively to the medical imaging domain.
	\item \textbf{Grad-CAM Compatibility:} The architecture's final convolutional
	layer (\texttt{block14\_sepconv2\_act}) produces spatially meaningful 2D
	feature maps that directly feed the Grad-CAM localization algorithm.
\end{itemize}

A custom classifier head was appended to the Xception backbone: a
\texttt{BatchNormalization} layer, followed by two \texttt{Dense} layers (256 and
128 neurons) with \texttt{Swish} activation, L2 regularization
($\lambda = 0.0016$), and \texttt{Dropout} (rates 0.5 and 0.2 respectively),
terminated by a 7-neuron \texttt{Softmax} output layer.

% ────────────────────────────────────────────────────────────
\subsection{Training Configuration \& Adaptive LR Strategy}
% ────────────────────────────────────────────────────────────
The model was compiled using the \textbf{AdamW} optimizer (Adam with decoupled
weight decay), which prevents weight decay from interfering with the gradient
update step and generally improves generalization~\cite{TensorFlow}. The model was trained using Categorical Cross-Entropy loss, a batch size of 16, and an initial learning rate of 0.001 over a maximum of 50 epochs.

A custom \textbf{Learning Rate Annealing (LRA)} callback was implemented to
dynamically adapt the learning rate. When training accuracy is below 90\%, the
callback monitors training accuracy. Once the 90\% threshold is surpassed, the
monitor switches to \textbf{validation loss}, preventing the model from
overfitting to the training distribution during the fine-tuning phase.

% ────────────────────────────────────────────────────────────
\subsection{Training Progress}
% ────────────────────────────────────────────────────────────
Training ran for \textbf{18 epochs} before the early stopping criterion was
triggered (3 consecutive learning rate adjustments without validation improvement).
The total training duration was \textbf{1 hour, 22 minutes, and 3.89 seconds},
with each epoch averaging approximately 271 seconds.



Figure~\ref{fig:sprint2_training_curves} shows the complete training and validation
loss and accuracy progression across all 18 epochs, clearly showing convergence and
the effect of adaptive learning rate reductions.

\begin{figure}[H]
	\centering
	\optionalimage[0.85\textwidth]{chap_4/images/graph.png}{Training and Validation Curves}
	\caption{\textit{Training and Validation Loss \& Accuracy Curves (18 Epochs)}}
	\label{fig:sprint2_training_curves}
\end{figure}

% ────────────────────────────────────────────────────────────
\subsection{Evaluation on the Test Set}
% ────────────────────────────────────────────────────────────
The final model was evaluated on the held-out test set of \textbf{1,002 images},
achieving an overall test accuracy of \textbf{83.83\%}. The detailed per-class
performance metrics are presented in Table~\ref{tab:classification_report}.

\begin{table}[H]
	\centering
	\renewcommand{\arraystretch}{1.5}
	\begin{tabular}{|l|c|c|c|c|}
		\hline
		\rowcolor[HTML]{1F3864}
		{\color{white}\textbf{Class}} &
		{\color{white}\textbf{Precision}} &
		{\color{white}\textbf{Recall}} &
		{\color{white}\textbf{F1-Score}} &
		{\color{white}\textbf{Test Samples}} \\
		\hline
		AKIEC & 0.71 & 0.69 & 0.70 &  32 \\
		\hline
		BCC   & 0.83 & 0.85 & 0.84 &  52 \\
		\hline
		BKL   & 0.64 & 0.91 & 0.75 & 110 \\
		\hline
		DF    & 0.67 & 0.36 & 0.47 &  11 \\
		\hline
		MEL   & 0.54 & 0.62 & 0.58 & 112 \\
		\hline
		NV    & 0.96 & 0.87 & 0.91 & 671 \\
		\hline
		VASC  & 0.93 & 1.00 & 0.97 &  14 \\
		\hline
		\rowcolor[HTML]{EEF6FF}
		\textbf{Weighted Avg} & \textbf{0.86} & \textbf{0.84} & \textbf{0.84} & \textbf{1,002} \\
		\hline
		\rowcolor[HTML]{EEF6FF}
		\textbf{Macro Avg}    & \textbf{0.75} & \textbf{0.76} & \textbf{0.75} & \textbf{1,002} \\
		\hline
	\end{tabular}
	\caption{\textit{Per-Class Classification Report on the 1,002-Image Test Set}}
	\label{tab:classification_report}
\end{table}

Key observations from the results include:
\begin{itemize}
	\item \textbf{NV} achieves the highest F1-score (0.91) due to its dominant
	representation in the dataset.
	\item \textbf{VASC} achieves perfect recall (1.00), demonstrating that the
	augmentation strategy successfully compensated for its very limited
	original sample count (114 training images).
	\item \textbf{MEL (Melanoma)} achieves a recall of 0.62. From a clinical
	safety perspective, this is the most critical metric: a missed melanoma
	is far more dangerous than a false positive. The platform therefore
	displays an explicit ``Consult a specialist'' advisory for all
	predictions with confidence below 85\%.
	\item \textbf{DF (Dermatofibroma)} has the lowest performance (F1 = 0.47)
	due to its very limited support (only 11 test samples), making
	statistical evaluation unreliable for this class.
\end{itemize}

Figure~\ref{fig:confusion_matrix} presents the confusion matrix, providing a
visual breakdown of correct predictions versus misclassifications across all
seven classes.

\begin{figure}[H]
	\centering
	\optionalimage[0.78\textwidth]{chap_4/images/confusion_matrix.png}{Confusion Matrix}
	\caption{\textit{Confusion Matrix on the 1,002-Image Test Set}}
	\label{fig:confusion_matrix}
\end{figure}

% ────────────────────────────────────────────────────────────
\subsection{Model Explainability: Grad-CAM}
% ────────────────────────────────────────────────────────────
To establish clinical trust, the platform does not operate as a ``black box.''
The \textbf{Gradient-weighted Class Activation Mapping (Grad-CAM)}
algorithm~\cite{Selvaraju2017} was implemented. The process involves: (1)~a
forward pass through the full Xception network; (2)~extraction of the spatial
gradient of the predicted class score with respect to the output of the final
convolutional layer (\texttt{block14\_sepconv2\_act}); (3)~computation of a
weighted sum of the feature maps using these gradients as importance weights;
(4)~bilinear upsampling to the original image resolution; and (5)~application
of a \texttt{COLORMAP\_JET} color overlay where red areas indicate the regions
most influential for the prediction.

% ────────────────────────────────────────────────────────────
\subsection{Inference Performance}
% ────────────────────────────────────────────────────────────
The deployed model is loaded once at Flask server startup (taking approximately 10 seconds) and held in memory. During active operation, the total API response time (end-to-end), including image preprocessing, Xception inference, and Grad-CAM generation, averages under 2 seconds.

% ============================================================
\section{Progress Tracking}
% ============================================================
Task execution was monitored via the Trello Kanban board. Figure~\ref{fig:trello_sprint2} captures the state of the board at the completion of the sprint, verifying that all AI integration and frontend scanning tasks were moved to the "Done" list.

\begin{figure}[H]
	\centering
	\optionalimage[0.40\textwidth]{chap_4/images/trello_sprint2.png}{Trello Sprint 2}
	\caption{\textit{Trello Board at the End of Sprint 2}}
	\label{fig:trello_sprint2}
\end{figure}

% ============================================================
\section{Conclusion}
% ============================================================
Sprint 2 successfully delivered the core diagnostic functionality of the DermaAI platform through a fine-tuned Xception model. By pairing high-accuracy predictions with transparent Grad-CAM visualizations, the system provides explainable and trustworthy assessments. With the scanning engine operational, the next chapter focuses on expanding user assistance through an AI chatbot.
